\chapter{Resultados --- Jaime (seguridad y backend)}
\label{cap:resultados-jaime}

Este cap\'itulo documenta, con evidencia verificable en el c\'odigo y en
\texttt{docs/mediciones/sec/}, el trabajo de seguridad, autenticaci\'on,
pruebas y cobertura desarrollado por Jaime durante la Tercera Entrega. La
auditor\'ia se organiz\'o sobre las categor\'ias aplicables de OWASP Top~10~\cite{owasp-top10}.

\section{Autenticaci\'on JWT por cookies}

El backend emite JWT firmados con HMAC-SHA256 (\texttt{jjwt}~0.12.6). Cada
token contiene \textbf{10 claims}: los \textbf{7 claims est\'andar}
registrados por RFC~7519~\cite{rfc7519} (\texttt{iss}, \texttt{sub}, \texttt{aud},
\texttt{iat}, \texttt{nbf}, \texttt{exp}, \texttt{jti}) y \textbf{3 claims
personalizados} (\texttt{email}, \texttt{rol}, \texttt{typ}). El tipo de
token (\texttt{access} o \texttt{refresh}) se distingue \'unicamente por el
valor del claim \texttt{typ}: ambos comparten el mismo conjunto de nombres
de claims, generados por el mismo m\'etodo (\texttt{buildToken}), y solo
difieren en su tiempo de vida y en ese valor (ADR-006).

Los tokens se transportan en cookies \texttt{access_token} y
\texttt{refresh_token}, ambas con los atributos
\texttt{HttpOnly}+\texttt{Secure}+\texttt{SameSite=Strict}
(\texttt{JwtCookieService}), con rutas diferenciadas
(\texttt{Path=/} para el access token, \texttt{Path=/api/auth} para el
refresh token). El frontend nunca almacena el JWT en \texttt{localStorage};
el navegador env\'ia las cookies autom\'aticamente
(\texttt{credentials: include}). El header \texttt{Authorization: Bearer}
sigue soport\'andose como mecanismo secundario en
\texttt{JwtAuthenticationFilter}, pero el flujo web principal de BIOPET no
lo utiliza.

% EVIDENCIA-JAIME-08
% Captura: atributos de cookies ocultando completamente sus valores.
% Ruta esperada: figuras/jaime/08-cookie-attributes.png
\evidencia{figuras/jaime/08-cookie-attributes.png}%
{Atributos reales de las cookies de sesi\'on (\texttt{HttpOnly}, \texttt{Secure}, \texttt{SameSite=Strict}) inspeccionados en el navegador, sin exponer el valor del token.}%
{fig:jaime-cookie-attributes}

\subsection{Logout y revocaci\'on en Redis}

El logout (\texttt{POST /api/auth/logout}) es idempotente: responde
\texttt{204} incluso sin cookies previas. Cuando hay un token v\'alido, se
extrae su \texttt{jti} y se marca como revocado en Redis
(\texttt{TokenBlacklistService}) con un TTL igual al tiempo de vida
restante del token. En cada solicitud protegida,
\texttt{JwtAuthenticationFilter} consulta esa lista negra \'unicamente para
tokens de tipo \texttt{access}, evitando consultas innecesarias a Redis
para refresh tokens.

\section{Autorizaci\'on por rol y por propietario}

La autorizaci\'on combina dos niveles: por \textbf{rol}
(\texttt{@PreAuthorize} en \texttt{MascotaController}, con
\texttt{ADMIN}/\texttt{VETERINARIO}/\texttt{AUXILIAR} habilitados para
escritura y los cuatro roles habilitados para lectura), y por
\textbf{propiedad} (\texttt{MascotaService.verificarPropiedad}, que lanza
\texttt{AccessDeniedException} cuando un \texttt{ROLE_DUENO} intenta
acceder a una mascota que no le pertenece). Las respuestas se diferencian
correctamente: \textbf{401} cuando no hay autenticaci\'on v\'alida
(\texttt{ProblemAuthenticationEntryPoint}), \textbf{403} cuando hay
autenticaci\'on pero no autorizaci\'on (\texttt{ProblemAccessDeniedHandler}).
Ambos casos producen un \texttt{ProblemDetail} real, sin exposici\'on de
\emph{stack trace}.

% EVIDENCIA-JAIME-05
% Captura: Postman con respuesta 401 ProblemDetail.
% Ruta esperada: figuras/jaime/05-postman-401.png
\evidencia{figuras/jaime/05-postman-401.png}%
{Respuesta 401 (\texttt{ProblemDetail}, \texttt{type=urn:biopet:error:unauthorized}) capturada desde la colecci\'on Postman de BIOPET.}%
{fig:jaime-postman-401}

% EVIDENCIA-JAIME-06
% Captura: Postman con respuesta 403 ProblemDetail.
% Ruta esperada: figuras/jaime/06-postman-403.png
\evidencia{figuras/jaime/06-postman-403.png}%
{Respuesta 403 (\texttt{ProblemDetail}, \texttt{type=urn:biopet:error:forbidden}) capturada desde la colecci\'on Postman de BIOPET.}%
{fig:jaime-postman-403}

\section{Rate limiting de login}

\texttt{LoginRateLimiterService} mantiene, en memoria
(\texttt{ConcurrentHashMap}) y por IP, un contador de intentos fallidos con
ventana y bloqueo de 15~minutos por defecto (\texttt{maxAttempts=6}): los
primeros 5~fallos consecutivos responden \texttt{401}, el 6.\textsuperscript{o}
responde \texttt{429} con cabecera \texttt{Retry-After} (entero positivo de
segundos). El contador se reinicia tras un login exitoso y es
completamente independiente entre IP distintas. Al ser en memoria y por
instancia, no es distribuido entre m\'ultiples r\'eplicas del backend
(limitaci\'on documentada desde la Fase~6).

% EVIDENCIA-JAIME-07
% Captura: sexto intento de login con 429 y Retry-After.
% Ruta esperada: figuras/jaime/07-rate-limit-429.png
\evidencia{figuras/jaime/07-rate-limit-429.png}%
{Sexto intento fallido de login consecutivo: respuesta 429 con cabecera \texttt{Retry-After}.}%
{fig:jaime-rate-limit-429}

\section{\texttt{ProblemDetail} uniforme}

Todos los errores de la API usan el formato RFC~7807~\cite{rfc7807}
(\texttt{application/problem+json}), con los campos \texttt{type},
\texttt{title}, \texttt{status}, \texttt{detail} e \texttt{instance}
(\texttt{GlobalExceptionHandler}, \texttt{ProblemDetailFactory},
\texttt{ProblemType}). El \texttt{detail} nunca incluye \emph{stack trace}
ni el valor bruto de un par\'ametro inv\'alido: para
\texttt{MethodArgumentTypeMismatchException}, el mensaje se construye
\'unicamente con el nombre del par\'ametro (\texttt{ex.getName()}), nunca
con el valor recibido.

\section{Auditor\'ia A09}

\texttt{AuthenticationAuditService} registra siete eventos estructurados en
una sola l\'inea de log (\texttt{AUTH\_AUDIT
timestamp=<UTC> event=<EVENTO> result=<RESULTADO> ip=<IP> subject=<SUJETO>}):
\texttt{LOGIN_SUCCESS}, \texttt{LOGIN_FAILURE}, \texttt{LOGIN_RATE_LIMITED},
\texttt{REFRESH_SUCCESS}, \texttt{REFRESH_FAILURE}, \texttt{LOGOUT_SUCCESS}
y \texttt{TOKEN_REVOKED}, con nivel \texttt{INFO} para los eventos de
resultado esperado y \texttt{WARN} para los fallidos o bloqueados. El
servicio sanitiza \texttt{ip}/\texttt{subject} (elimina caracteres de
control, trunca a 200 caracteres, sustituye valores ausentes por
\texttt{"unknown"}) y \textbf{nunca} recibe como argumento una contrase\~na,
un JWT completo, una cookie ni el encabezado \texttt{Authorization}: la
firma de sus m\'etodos p\'ublicos solo acepta \texttt{ip} y \texttt{subject}.

% EVIDENCIA-JAIME-09
% Captura: log AUTH_AUDIT sin secretos.
% Ruta esperada: figuras/jaime/09-auth-audit.png
\evidencia{figuras/jaime/09-auth-audit.png}%
{L\'inea real de log \texttt{AUTH\_AUDIT} generada durante una prueba de login, sin contrase\~nas ni tokens.}%
{fig:jaime-auth-audit}

\section{Cabeceras de seguridad y TLS}

\texttt{SecurityConfig} (Spring Security~\cite{spring-security-docs}) fija cinco cabeceras en todas las respuestas
(incluidas las p\'ublicas): \texttt{X-Frame-Options: DENY},
\texttt{X-Content-Type-Options: nosniff}, \texttt{Content-Security-Policy:
default-src 'self'; frame-ancestors 'none'; object-src 'none'},
\texttt{Referrer-Policy: no-referrer}, y
\texttt{Strict-Transport-Security} \'unicamente sobre HTTPS
(\texttt{max-age=31536000 ; includeSubDomains ; preload}; ausente sobre
HTTP, comportamiento correcto y verificado).

% EVIDENCIA-JAIME-04
% Captura: curl HTTPS con HSTS y cabeceras de seguridad.
% Ruta esperada: figuras/jaime/04-security-headers.png
\evidencia{figuras/jaime/04-security-headers.png}%
{Cabeceras de seguridad reales capturadas con \texttt{curl.exe -i} contra \texttt{https://localhost:8443/actuator/health}.}%
{fig:jaime-security-headers}

El perfil \texttt{tls} activa un conector HTTPS en el puerto~8443
(\texttt{TLSv1.3} \'unicamente) junto al conector HTTP interno en el
puerto~8080. Verificado en vivo con \texttt{openssl s\_client}: protocolo
\textbf{TLSv1.3}, cifrado \textbf{\texttt{TLS\_AES\_256\_GCM\_SHA384}}
(suite AEAD), certificado autofirmado (\texttt{verify error:num=18:
self-signed certificate}) con SAN \texttt{DNS:localhost,
IP Address:127.0.0.1} --- exclusivamente acad\'emico, nunca v\'alido para
producci\'on.

% EVIDENCIA-JAIME-03
% Captura: OpenSSL mostrando TLSv1.3 y TLS_AES_256_GCM_SHA384.
% Ruta esperada: figuras/jaime/03-tls-openssl.png
\evidencia{figuras/jaime/03-tls-openssl.png}%
{Salida de \texttt{openssl s\_client -connect localhost:8443 -tls1\_3}: protocolo TLSv1.3 y cifrado \texttt{TLS\_AES\_256\_GCM\_SHA384} negociados.}%
{fig:jaime-tls-openssl}

\section{Protecci\'on estructural contra inyecci\'on (A03)}

Todos los m\'etodos de \texttt{UsuarioRepository} y \texttt{MascotaRepository}
son consultas derivadas de Spring Data, con par\'ametros enlazados. La
\'unica \texttt{@Query} nativa del proyecto
(\texttt{fn\_resumen\_mascotas\_por\_especie}) usa un par\'ametro
\texttt{:duenioId} enlazado y fuertemente tipado como \texttt{Long}. Los
payloads de inyecci\'on probados (\texttt{1 OR 1=1},
\texttt{1; DROP TABLE usuarios; --}, \texttt{\%' UNION SELECT NULL --})
nunca llegan a ejecutarse como SQL: son rechazados en la capa de conversi\'on
de Spring MVC (\texttt{MethodArgumentTypeMismatchException}, 400) antes de
alcanzar el repositorio, verificado en
\texttt{SqlInjectionSecurityTest} (incluye la comprobaci\'on de que
\texttt{usuarioRepository.count()} y \texttt{mascotaRepository.count()}
permanecen id\'enticos antes y despu\'es de cada payload).

\section{Postman}

La colecci\'on \texttt{docs/postman/BIOPET.postman\_collection.json}
(Schema~v2.1.0, 40~requests en 5~carpetas) cubre estado del servicio,
autenticaci\'on (incluida la secuencia completa de rate limiting),
mascotas, resumen por especie, y un cat\'alogo de errores controlados
(422, 400, 404, 401, 403, 429). Depende exclusivamente del \emph{cookie
jar} autom\'atico de Postman: ning\'un script guarda un JWT o el valor de
una cookie en una variable, y ning\'un request agrega
\texttt{Authorization: Bearer} al flujo principal.

\section{Cobertura JaCoCo y resultado de pruebas}

JaCoCo~\cite{jacoco-docs} es la herramienta usada para medir y verificar
autom\'aticamente esta cobertura.

\begin{table}[H]
  \centering
  \caption{Cobertura medida por \texttt{jacoco:check} (\texttt{Backend/target/site/jacoco/jacoco.xml}, verificado en esta fase).}
  \label{tab:jacoco-jaime}
  \begin{tabular}{@{}lrrrr@{}}
    \toprule
    Contador & Cubierto & No cubierto & Total & Porcentaje \\
    \midrule
    LINE & 534 & 23 & 557 & 95{,}87\,\% \\
    BRANCH & 103 & 31 & 134 & 76{,}87\,\% \\
    COMPLEXITY & 172 & 43 & 215 & 80{,}00\,\% \\
    \bottomrule
  \end{tabular}
\end{table}

Umbral autom\'atico configurado: \(\geq\) 60\,\% para los tres contadores
anteriores (regla \texttt{BUNDLE} de \texttt{jacoco-maven-plugin}). Las 109
pruebas de la suite completa terminaron con 0~fallos, 0~errores y
0~omitidas (\texttt{BUILD SUCCESS}).

% EVIDENCIA-JAIME-01
% Captura: mvn clean verify con 109 pruebas y BUILD SUCCESS.
% Ruta esperada: figuras/jaime/01-maven-verify.png
\evidencia{figuras/jaime/01-maven-verify.png}%
{Salida real de \texttt{mvn clean verify}: 109 pruebas, 0 fallos, 0 errores, 0 omitidas, \texttt{BUILD SUCCESS}.}%
{fig:jaime-maven-verify}

% EVIDENCIA-JAIME-02
% Captura: resumen HTML de JaCoCo.
% Ruta esperada: figuras/jaime/02-jacoco-resumen.png
\evidencia{figuras/jaime/02-jacoco-resumen.png}%
{Resumen navegable de JaCoCo (\texttt{Backend/target/site/jacoco/index.html}) con la cobertura por paquete.}%
{fig:jaime-jacoco-resumen}

\section{C4 Nivel 3 y ADR-006}

El diagrama de componentes del backend (cap\'itulo~\ref{cap:arquitectura-c4})
y ADR-006 (\emph{Decisiones de autenticaci\'on y seguridad}) documentan
formalmente las decisiones descritas en este cap\'itulo: cookies frente a
almacenamiento en el cliente, revocaci\'on v\'ia Redis, rate limiting en
memoria, y el formato \texttt{ProblemDetail} para todos los errores.

% EVIDENCIA-JAIME-10
% Captura: C4 Nivel 3 renderizado.
% Ruta esperada: figuras/jaime/10-c4-nivel-3.png
\evidencia{figuras/jaime/10-c4-nivel-3.png}%
{Diagrama C4 Nivel 3 (componentes del backend) renderizado, referenciado tambi\'en en la secci\'on~\ref{sec:c4-nivel-3}.}%
{fig:jaime-c4-nivel-3-resultados}

\section{Limitaciones reales de este bloque}

\begin{itemize}[nosep]
  \item El certificado TLS es autofirmado y exclusivamente acad\'emico/local;
    no representa una cadena de confianza v\'alida para producci\'on.
  \item El rate limiting de login es en memoria y por instancia: no es
    distribuido entre m\'ultiples r\'eplicas del backend.
  \item Los logs \texttt{AUTH_AUDIT} se escriben \'unicamente en el log
    local del proceso; no existe integraci\'on con un SIEM centralizado ni
    alertas autom\'aticas.
  \item La revocaci\'on de tokens depende de la disponibilidad de Redis; si
    Redis no est\'a disponible, debe tratarse como error de infraestructura,
    no como sesi\'on v\'alida (ADR-003).
  \item Toda la evidencia de este cap\'itulo se recolect\'o contra
    \textbf{una sola instancia} del backend en un entorno local; no se
    evalu\'o comportamiento bajo m\'ultiples r\'eplicas ni balanceo de
    carga.
\end{itemize}
